%%%%%%%%%%%%%%%%%%%%%%%%%%%%%%%%%%%%%%%%%%%%%%%%%%%%%%%%%%%%
% 8.5 CONFIDENCE INTERVALS
% The Composition of Advanced Mathematics
%%%%%%%%%%%%%%%%%%%%%%%%%%%%%%%%%%%%%%%%%%%%%%%%%%%%%%%%%%%%

\section{Confidence Intervals}
\label{sec:confidence-intervals}

\prerequisite{
Sampling distributions, statistical estimation, normal distributions,
Student's \(t\)-distribution, chi-square and \(F\)-distributions,
standard errors, point estimation, central limit theorems, and basic
probability.
}

\learningobjectives{
By the end of this section, the reader should be able to:
\begin{itemize}
    \item explain the repeated-sampling meaning of a confidence interval;
    \item distinguish confidence level from probability assigned to a
          fixed parameter;
    \item define margin of error;
    \item construct confidence intervals from pivotal quantities;
    \item construct \(z\)-intervals for a population mean with known
          standard deviation;
    \item construct one-sample \(t\)-intervals for a population mean with
          unknown standard deviation;
    \item construct large-sample confidence intervals for population
          proportions;
    \item understand Wilson and exact binomial intervals conceptually;
    \item construct intervals for differences of two means;
    \item distinguish pooled and Welch two-sample intervals;
    \item construct paired confidence intervals;
    \item construct intervals for differences of proportions;
    \item construct normal-theory confidence intervals for a variance;
    \item construct confidence intervals for ratios of variances;
    \item understand likelihood-based confidence intervals;
    \item understand bootstrap confidence intervals;
    \item distinguish percentile and basic bootstrap intervals;
    \item understand transformation-based intervals;
    \item interpret confidence intervals correctly;
    \item distinguish statistical precision from practical importance;
    \item determine approximate sample sizes for desired margins of error;
    \item understand one-sided confidence bounds;
    \item connect confidence intervals with hypothesis tests and
          estimation.
\end{itemize}
}

%%%%%%%%%%%%%%%%%%%%%%%%%%%%%%%%%%%%%%%%%%%%%%%%%%%%%%%%%%%%
\subsection{Why Interval Estimation?}
\label{subsec:why-interval-estimation}
%%%%%%%%%%%%%%%%%%%%%%%%%%%%%%%%%%%%%%%%%%%%%%%%%%%%%%%%%%%%

A point estimate gives a single numerical guess for an unknown
parameter.

For example,

\[
\widehat{\mu}
=
\overline{X}.
\]

However, two samples from the same population generally produce
different sample means.

Thus a point estimate alone does not communicate the amount of sampling
uncertainty.

A confidence interval supplements the point estimate with a range of
values compatible with the observed data under a specified sampling
procedure.

%%%%%%%%%%%%%%%%%%%%%%%%%%%%%%%%%%%%%%%%%%%%%%%%%%%%%%%%%%%%
\subsection{General Form of a Confidence Interval}
\label{subsec:general-confidence-interval}
%%%%%%%%%%%%%%%%%%%%%%%%%%%%%%%%%%%%%%%%%%%%%%%%%%%%%%%%%%%%

Many confidence intervals have the form

\[
\boxed{
\text{estimate}
\pm
\text{critical value}
\times
\text{standard error}.
}
\]

The quantity

\[
\boxed{
\text{critical value}
\times
\text{standard error}
}
\]

is called the margin of error.

Thus,

\[
\boxed{
\text{CI}
=
\text{estimate}
\pm
\text{margin of error}.
}
\]

%%%%%%%%%%%%%%%%%%%%%%%%%%%%%%%%%%%%%%%%%%%%%%%%%%%%%%%%%%%%
\subsection{Confidence Coefficient}
\label{subsec:confidence-coefficient}
%%%%%%%%%%%%%%%%%%%%%%%%%%%%%%%%%%%%%%%%%%%%%%%%%%%%%%%%%%%%

A confidence level is usually written

\[
\boxed{
1-\alpha,
}
\]

where

\[
0<\alpha<1.
\]

Common choices include

\[
0.90,
\qquad
0.95,
\qquad
0.99.
\]

Thus a \(95\%\) confidence procedure uses

\[
\alpha=0.05.
\]

%%%%%%%%%%%%%%%%%%%%%%%%%%%%%%%%%%%%%%%%%%%%%%%%%%%%%%%%%%%%
\subsection{Repeated-Sampling Interpretation}
\label{subsec:repeated-sampling-confidence}
%%%%%%%%%%%%%%%%%%%%%%%%%%%%%%%%%%%%%%%%%%%%%%%%%%%%%%%%%%%%

Suppose a confidence interval procedure produces random endpoints

\[
L(X_1,\ldots,X_n)
\]

and

\[
U(X_1,\ldots,X_n).
\]

A \(100(1-\alpha)\%\) confidence procedure satisfies

\[
\boxed{
P_\theta
\left(
L(X)\leq\theta\leq U(X)
\right)
=
1-\alpha
}
\]

for the specified model, exactly or approximately.

The probability refers to the randomness of the interval before the
sample is observed.

%%%%%%%%%%%%%%%%%%%%%%%%%%%%%%%%%%%%%%%%%%%%%%%%%%%%%%%%%%%%
\subsection{Correct Interpretation After Observing the Data}
\label{subsec:correct-confidence-interpretation}
%%%%%%%%%%%%%%%%%%%%%%%%%%%%%%%%%%%%%%%%%%%%%%%%%%%%%%%%%%%%

After the data are observed, the interval endpoints become fixed.

The parameter is also treated as fixed in frequentist inference.

Therefore one should not usually say

\[
\text{``There is a \(95\%\) probability that \(\theta\) lies in this
interval.''}
\]

Instead, say:

\[
\boxed{
\text{This interval was produced by a procedure that captures the true
parameter in \(95\%\) of repeated samples under the model.}
}
\]

%%%%%%%%%%%%%%%%%%%%%%%%%%%%%%%%%%%%%%%%%%%%%%%%%%%%%%%%%%%%
\subsection{Confidence Level and Interval Width}
\label{subsec:confidence-level-width}
%%%%%%%%%%%%%%%%%%%%%%%%%%%%%%%%%%%%%%%%%%%%%%%%%%%%%%%%%%%%

Increasing the confidence level requires a larger critical value.

Therefore, with all else fixed,

\[
\boxed{
\text{higher confidence}
\Longrightarrow
\text{wider interval}.
}
\]

Conversely,

\[
\boxed{
\text{lower confidence}
\Longrightarrow
\text{narrower interval}.
}
\]

There is a tradeoff between confidence and precision.

%%%%%%%%%%%%%%%%%%%%%%%%%%%%%%%%%%%%%%%%%%%%%%%%%%%%%%%%%%%%
\subsection{Critical Values from the Standard Normal}
\label{subsec:z-critical-values}
%%%%%%%%%%%%%%%%%%%%%%%%%%%%%%%%%%%%%%%%%%%%%%%%%%%%%%%%%%%%

Let

\[
Z\sim N(0,1).
\]

For a two-sided confidence level \(1-\alpha\), define

\[
z_{\alpha/2}
\]

or equivalently \(z^*\) by

\[
\boxed{
P
(
-Z^*
\leq
Z
\leq
Z^*
)
=
1-\alpha.
}
\]

Equivalently,

\[
\boxed{
P(Z>Z^*)
=
\frac{\alpha}{2}.
}
\]

Common values are approximately

\[
\begin{array}{c|c}
\text{Confidence Level} & z^*\\
\hline
90\% & 1.645\\
95\% & 1.960\\
99\% & 2.576
\end{array}
\]

%%%%%%%%%%%%%%%%%%%%%%%%%%%%%%%%%%%%%%%%%%%%%%%%%%%%%%%%%%%%
\subsection{Confidence Interval for a Mean with Known \(\sigma\)}
\label{subsec:ci-mean-known-sigma}
%%%%%%%%%%%%%%%%%%%%%%%%%%%%%%%%%%%%%%%%%%%%%%%%%%%%%%%%%%%%

Suppose

\[
X_1,\ldots,X_n
\]

are sampled from a normal population with mean \(\mu\) and known
standard deviation \(\sigma\).

Then

\[
\boxed{
Z
=
\frac{
\overline{X}-\mu
}{
\sigma/\sqrt{n}
}
\sim
N(0,1).
}
\]

Thus,

\[
P
\left(
-z^*
\leq
\frac{
\overline{X}-\mu
}{
\sigma/\sqrt{n}
}
\leq
z^*
\right)
=
1-\alpha.
\]

Solving for \(\mu\),

\[
\boxed{
\overline{X}
-
z^*
\frac{\sigma}{\sqrt{n}}
\leq
\mu
\leq
\overline{X}
+
z^*
\frac{\sigma}{\sqrt{n}}.
}
\]

Therefore the confidence interval is

\[
\boxed{
\overline{X}
\pm
z^*
\frac{\sigma}{\sqrt{n}}.
}
\]

%%%%%%%%%%%%%%%%%%%%%%%%%%%%%%%%%%%%%%%%%%%%%%%%%%%%%%%%%%%%
\subsection{Worked Example: Known-\(\sigma\) Mean Interval}
\label{subsec:worked-known-sigma-ci}
%%%%%%%%%%%%%%%%%%%%%%%%%%%%%%%%%%%%%%%%%%%%%%%%%%%%%%%%%%%%

\begin{workedexamplebox}{A \(95\%\) Confidence Interval for a Mean}

Suppose

\[
\overline{x}=72,
\qquad
\sigma=10,
\qquad
n=100.
\]

For \(95\%\) confidence,

\[
z^*=1.96.
\]

The standard error is

\[
\frac{10}{\sqrt{100}}
=
1.
\]

The margin of error is

\[
1.96(1)
=
1.96.
\]

Therefore,

\begin{answerbox}
\[
\boxed{
72\pm1.96
}
\]

or

\[
\boxed{
(70.04,\ 73.96).
}
\]
\end{answerbox}

\end{workedexamplebox}

%%%%%%%%%%%%%%%%%%%%%%%%%%%%%%%%%%%%%%%%%%%%%%%%%%%%%%%%%%%%
\subsection{Margin of Error for a Mean}
\label{subsec:margin-error-mean}
%%%%%%%%%%%%%%%%%%%%%%%%%%%%%%%%%%%%%%%%%%%%%%%%%%%%%%%%%%%%

For a known-\(\sigma\) mean interval,

\[
\boxed{
E
=
z^*
\frac{\sigma}{\sqrt{n}}.
}
\]

The margin of error decreases when:

\[
\boxed{
n
\text{ increases},
}
\]

or

\[
\boxed{
\sigma
\text{ decreases}.
}
\]

It increases when the desired confidence level increases.

%%%%%%%%%%%%%%%%%%%%%%%%%%%%%%%%%%%%%%%%%%%%%%%%%%%%%%%%%%%%
\subsection{Sample Size for a Desired Mean Margin of Error}
\label{subsec:sample-size-mean-ci}
%%%%%%%%%%%%%%%%%%%%%%%%%%%%%%%%%%%%%%%%%%%%%%%%%%%%%%%%%%%%

Suppose the desired margin of error is at most \(E\).

From

\[
E
=
z^*
\frac{\sigma}{\sqrt{n}},
\]

solve for \(n\):

\[
\sqrt{n}
=
\frac{
z^*\sigma
}{E},
\]

so

\[
\boxed{
n
=
\left(
\frac{
z^*\sigma
}{E}
\right)^2.
}
\]

Because sample size must be an integer, round upward.

%%%%%%%%%%%%%%%%%%%%%%%%%%%%%%%%%%%%%%%%%%%%%%%%%%%%%%%%%%%%
\subsection{Worked Example: Planning a Mean Sample Size}
\label{subsec:worked-sample-size-mean}
%%%%%%%%%%%%%%%%%%%%%%%%%%%%%%%%%%%%%%%%%%%%%%%%%%%%%%%%%%%%

\begin{workedexamplebox}{Desired Margin of Error}

Suppose a \(95\%\) interval is desired with margin of error at most

\[
E=2.
\]

Assume

\[
\sigma=12.
\]

Then

\[
n
=
\left(
\frac{
1.96(12)
}{2}
\right)^2.
\]

Thus,

\[
n
=
(11.76)^2
\approx
138.30.
\]

Round upward:

\begin{answerbox}
\[
\boxed{
n=139.
}
\]
\end{answerbox}

\end{workedexamplebox}

%%%%%%%%%%%%%%%%%%%%%%%%%%%%%%%%%%%%%%%%%%%%%%%%%%%%%%%%%%%%
\subsection{Unknown Population Standard Deviation}
\label{subsec:unknown-sigma-ci}
%%%%%%%%%%%%%%%%%%%%%%%%%%%%%%%%%%%%%%%%%%%%%%%%%%%%%%%%%%%%

In most real applications,

\[
\sigma
\]

is unknown.

Replacing \(\sigma\) by the sample standard deviation \(S\) introduces
additional uncertainty.

For normal samples,

\[
\boxed{
T
=
\frac{
\overline{X}-\mu
}{
S/\sqrt{n}
}
\sim
t_{n-1}.
}
\]

This leads to the one-sample \(t\)-interval.

%%%%%%%%%%%%%%%%%%%%%%%%%%%%%%%%%%%%%%%%%%%%%%%%%%%%%%%%%%%%
\subsection{One-Sample \(t\)-Interval}
\label{subsec:one-sample-t-interval}
%%%%%%%%%%%%%%%%%%%%%%%%%%%%%%%%%%%%%%%%%%%%%%%%%%%%%%%%%%%%

For a normal population with unknown standard deviation,

\[
\boxed{
\mu
\in
\overline{X}
\pm
t^*_{n-1}
\frac{S}{\sqrt{n}},
}
\]

where \(t^*_{n-1}\) is the appropriate critical value from a
\(t\)-distribution with

\[
n-1
\]

degrees of freedom.

%%%%%%%%%%%%%%%%%%%%%%%%%%%%%%%%%%%%%%%%%%%%%%%%%%%%%%%%%%%%
\subsection{\(t\)-Critical Values}
\label{subsec:t-critical-values}
%%%%%%%%%%%%%%%%%%%%%%%%%%%%%%%%%%%%%%%%%%%%%%%%%%%%%%%%%%%%

For fixed confidence level, \(t\)-critical values depend on degrees of
freedom.

For small degrees of freedom,

\[
t^*
\]

is larger than the corresponding standard-normal critical value.

As

\[
n\to\infty,
\]

\[
\boxed{
t^*_{n-1}
\to
z^*.
}
\]

This reflects decreasing uncertainty from estimating \(\sigma\).

%%%%%%%%%%%%%%%%%%%%%%%%%%%%%%%%%%%%%%%%%%%%%%%%%%%%%%%%%%%%
\subsection{Worked Example: One-Sample \(t\)-Interval}
\label{subsec:worked-one-sample-t-ci}
%%%%%%%%%%%%%%%%%%%%%%%%%%%%%%%%%%%%%%%%%%%%%%%%%%%%%%%%%%%%

\begin{workedexamplebox}{Mean with Unknown Standard Deviation}

Suppose

\[
n=25,
\qquad
\overline{x}=50,
\qquad
s=10.
\]

For a \(95\%\) interval with

\[
df=24,
\]

suppose

\[
t^*
\approx2.064.
\]

The standard error is

\[
\frac{10}{5}=2.
\]

Thus the margin of error is

\[
2.064(2)
=
4.128.
\]

Therefore,

\begin{answerbox}
\[
\boxed{
50\pm4.128
}
\]

or approximately

\[
\boxed{
(45.872,\ 54.128).
}
\]
\end{answerbox}

\end{workedexamplebox}

%%%%%%%%%%%%%%%%%%%%%%%%%%%%%%%%%%%%%%%%%%%%%%%%%%%%%%%%%%%%
\subsection{Conditions for the One-Sample \(t\)-Interval}
\label{subsec:t-interval-conditions}
%%%%%%%%%%%%%%%%%%%%%%%%%%%%%%%%%%%%%%%%%%%%%%%%%%%%%%%%%%%%

The exact \(t\)-interval is derived under normal sampling.

For larger samples, the procedure is often approximately valid under
weaker conditions because of asymptotic normality and consistency of the
sample variance.

However, severe skewness, extreme outliers, or heavy tails may require
larger sample sizes or alternative methods.

The adequacy of the interval should be judged from both sample size and
distributional behavior.

%%%%%%%%%%%%%%%%%%%%%%%%%%%%%%%%%%%%%%%%%%%%%%%%%%%%%%%%%%%%
\subsection{Confidence Interval for a Population Proportion}
\label{subsec:ci-population-proportion}
%%%%%%%%%%%%%%%%%%%%%%%%%%%%%%%%%%%%%%%%%%%%%%%%%%%%%%%%%%%%

Suppose

\[
X\sim\operatorname{Bin}(n,p)
\]

and

\[
\widehat{p}
=
\frac Xn.
\]

For sufficiently large \(n\),

\[
\widehat{p}
\approx
N
\left(
p,
\frac{
p(1-p)
}{n}
\right).
\]

Replacing \(p\) in the standard error by \(\widehat{p}\) gives the
Wald interval:

\[
\boxed{
\widehat{p}
\pm
z^*
\sqrt{
\frac{
\widehat{p}(1-\widehat{p})
}{n}
}.
}
\]

%%%%%%%%%%%%%%%%%%%%%%%%%%%%%%%%%%%%%%%%%%%%%%%%%%%%%%%%%%%%
\subsection{Worked Example: Proportion Confidence Interval}
\label{subsec:worked-proportion-ci}
%%%%%%%%%%%%%%%%%%%%%%%%%%%%%%%%%%%%%%%%%%%%%%%%%%%%%%%%%%%%

\begin{workedexamplebox}{Estimating a Population Proportion}

Suppose

\[
n=400
\]

and

\[
x=240.
\]

Then

\[
\widehat{p}
=
\frac{240}{400}
=
0.60.
\]

The estimated standard error is

\[
\sqrt{
\frac{
(0.60)(0.40)
}{400}
}
=
\sqrt{0.0006}.
\]

For \(95\%\) confidence,

\[
z^*=1.96.
\]

Therefore,

\begin{answerbox}
\[
\boxed{
0.60
\pm
1.96
\sqrt{0.0006}.
}
\]
\end{answerbox}

\end{workedexamplebox}

%%%%%%%%%%%%%%%%%%%%%%%%%%%%%%%%%%%%%%%%%%%%%%%%%%%%%%%%%%%%
\subsection{Limitations of the Wald Proportion Interval}
\label{subsec:wald-proportion-limitations}
%%%%%%%%%%%%%%%%%%%%%%%%%%%%%%%%%%%%%%%%%%%%%%%%%%%%%%%%%%%%

The simple Wald interval is easy to compute but can perform poorly when:

\[
\boxed{
n
\text{ is small},
}
\]

or

\[
\boxed{
p
\text{ is near }0\text{ or }1.
}
\]

Its coverage probability may differ substantially from the nominal
confidence level.

Better alternatives include:

\[
\boxed{
\text{Wilson intervals},
}
\]

\[
\boxed{
\text{Agresti--Coull-type intervals},
}
\]

and

\[
\boxed{
\text{exact binomial intervals}.
}
\]

%%%%%%%%%%%%%%%%%%%%%%%%%%%%%%%%%%%%%%%%%%%%%%%%%%%%%%%%%%%%
\subsection{Wilson Score Interval}
\label{subsec:wilson-score-interval}
%%%%%%%%%%%%%%%%%%%%%%%%%%%%%%%%%%%%%%%%%%%%%%%%%%%%%%%%%%%%

The Wilson interval is obtained by inverting a score test.

For a two-sided standard-normal critical value \(z^*\), its center is

\[
\boxed{
\frac{
\widehat{p}
+
\frac{(z^*)^2}{2n}
}{
1+\frac{(z^*)^2}{n}
}.
}
\]

Its half-width is

\[
\boxed{
\frac{
z^*
}{
1+\frac{(z^*)^2}{n}
}
\sqrt{
\frac{
\widehat{p}(1-\widehat{p})
}{n}
+
\frac{
(z^*)^2
}{
4n^2
}
}.
}
\]

The Wilson interval often has much better coverage than the simple Wald
interval.

%%%%%%%%%%%%%%%%%%%%%%%%%%%%%%%%%%%%%%%%%%%%%%%%%%%%%%%%%%%%
\subsection{Exact Binomial Interval Preview}
\label{subsec:exact-binomial-ci-preview}
%%%%%%%%%%%%%%%%%%%%%%%%%%%%%%%%%%%%%%%%%%%%%%%%%%%%%%%%%%%%

An exact binomial interval can be obtained by inverting exact binomial
tail probabilities.

One common construction is the Clopper--Pearson interval.

Its coverage is at least the nominal level under the binomial model, but
it may be conservative.

That means the actual coverage probability can exceed the stated level.

%%%%%%%%%%%%%%%%%%%%%%%%%%%%%%%%%%%%%%%%%%%%%%%%%%%%%%%%%%%%
\subsection{Sample Size for a Proportion}
\label{subsec:sample-size-proportion}
%%%%%%%%%%%%%%%%%%%%%%%%%%%%%%%%%%%%%%%%%%%%%%%%%%%%%%%%%%%%

For the approximate Wald margin of error

\[
E
=
z^*
\sqrt{
\frac{
p(1-p)
}{n}
},
\]

solve for \(n\):

\[
\boxed{
n
=
\frac{
(z^*)^2p(1-p)
}{
E^2
}.
}
\]

If a planning value of \(p\) is unavailable, use

\[
\boxed{
p=0.5,
}
\]

because

\[
p(1-p)
\]

is maximized at \(p=0.5\).

%%%%%%%%%%%%%%%%%%%%%%%%%%%%%%%%%%%%%%%%%%%%%%%%%%%%%%%%%%%%
\subsection{Worked Example: Sample Size for a Proportion}
\label{subsec:worked-proportion-sample-size}
%%%%%%%%%%%%%%%%%%%%%%%%%%%%%%%%%%%%%%%%%%%%%%%%%%%%%%%%%%%%

\begin{workedexamplebox}{Planning Without a Prior Proportion Estimate}

Suppose a \(95\%\) interval is desired with margin of error at most

\[
E=0.03.
\]

Use the conservative value

\[
p=0.5.
\]

Then

\[
n
=
\frac{
(1.96)^2(0.5)(0.5)
}{
(0.03)^2
}.
\]

Thus,

\[
n
\approx
1067.11.
\]

Round upward:

\begin{answerbox}
\[
\boxed{
n=1068.
}
\]
\end{answerbox}

\end{workedexamplebox}

%%%%%%%%%%%%%%%%%%%%%%%%%%%%%%%%%%%%%%%%%%%%%%%%%%%%%%%%%%%%
\subsection{Confidence Interval for Difference of Two Means}
\label{subsec:ci-difference-two-means}
%%%%%%%%%%%%%%%%%%%%%%%%%%%%%%%%%%%%%%%%%%%%%%%%%%%%%%%%%%%%

Suppose independent populations have means

\[
\mu_1
\]

and

\[
\mu_2.
\]

The parameter of interest is

\[
\boxed{
\mu_1-\mu_2.
}
\]

The point estimator is

\[
\boxed{
\overline{X}_1-\overline{X}_2.
}
\]

%%%%%%%%%%%%%%%%%%%%%%%%%%%%%%%%%%%%%%%%%%%%%%%%%%%%%%%%%%%%
\subsection{Known-Variance Two-Mean Interval}
\label{subsec:known-variance-two-mean-ci}
%%%%%%%%%%%%%%%%%%%%%%%%%%%%%%%%%%%%%%%%%%%%%%%%%%%%%%%%%%%%

If the population variances are known,

\[
\boxed{
\operatorname{SE}
(
\overline{X}_1-\overline{X}_2
)
=
\sqrt{
\frac{
\sigma_1^2
}{
n_1
}
+
\frac{
\sigma_2^2
}{
n_2
}
}.
}
\]

Thus the confidence interval is

\[
\boxed{
(\overline{X}_1-\overline{X}_2)
\pm
z^*
\sqrt{
\frac{
\sigma_1^2
}{
n_1
}
+
\frac{
\sigma_2^2
}{
n_2
}
}.
}
\]

%%%%%%%%%%%%%%%%%%%%%%%%%%%%%%%%%%%%%%%%%%%%%%%%%%%%%%%%%%%%
\subsection{Welch Two-Sample \(t\)-Interval}
\label{subsec:welch-two-sample-ci}
%%%%%%%%%%%%%%%%%%%%%%%%%%%%%%%%%%%%%%%%%%%%%%%%%%%%%%%%%%%%

When the variances are unknown and not assumed equal, use

\[
\boxed{
(\overline{X}_1-\overline{X}_2)
\pm
t^*
\sqrt{
\frac{
S_1^2
}{
n_1
}
+
\frac{
S_2^2
}{
n_2
}
}.
}
\]

The degrees of freedom are approximated using the
Welch--Satterthwaite formula:

\[
\boxed{
\nu
\approx
\frac{
\left(
S_1^2/n_1
+
S_2^2/n_2
\right)^2
}{
\frac{
(S_1^2/n_1)^2
}{
n_1-1
}
+
\frac{
(S_2^2/n_2)^2
}{
n_2-1
}
}.
}
\]

%%%%%%%%%%%%%%%%%%%%%%%%%%%%%%%%%%%%%%%%%%%%%%%%%%%%%%%%%%%%
\subsection{Worked Example: Welch Interval}
\label{subsec:worked-welch-ci}
%%%%%%%%%%%%%%%%%%%%%%%%%%%%%%%%%%%%%%%%%%%%%%%%%%%%%%%%%%%%

\begin{workedexamplebox}{Difference Between Two Independent Means}

Suppose

\[
\overline{x}_1=80,
\qquad
s_1=12,
\qquad
n_1=40,
\]

and

\[
\overline{x}_2=74,
\qquad
s_2=10,
\qquad
n_2=35.
\]

The point estimate is

\[
80-74=6.
\]

The estimated standard error is

\[
\sqrt{
\frac{12^2}{40}
+
\frac{10^2}{35}
}.
\]

Thus the interval has the form

\begin{answerbox}
\[
\boxed{
6
\pm
t^*
\sqrt{
\frac{144}{40}
+
\frac{100}{35}
},
}
\]

with \(t^*\) based on the Welch degrees of freedom.
\end{answerbox}

\end{workedexamplebox}

%%%%%%%%%%%%%%%%%%%%%%%%%%%%%%%%%%%%%%%%%%%%%%%%%%%%%%%%%%%%
\subsection{Pooled Two-Sample \(t\)-Interval}
\label{subsec:pooled-two-sample-ci}
%%%%%%%%%%%%%%%%%%%%%%%%%%%%%%%%%%%%%%%%%%%%%%%%%%%%%%%%%%%%

If two independent normal populations are assumed to have a common
variance,

\[
\sigma_1^2
=
\sigma_2^2
=
\sigma^2,
\]

use the pooled variance estimator

\[
\boxed{
S_p^2
=
\frac{
(n_1-1)S_1^2
+
(n_2-1)S_2^2
}{
n_1+n_2-2
}.
}
\]

Then

\[
\boxed{
\operatorname{SE}
=
S_p
\sqrt{
\frac1{n_1}
+
\frac1{n_2}
}.
}
\]

The confidence interval is

\[
\boxed{
(\overline{X}_1-\overline{X}_2)
\pm
t^*_{n_1+n_2-2}
S_p
\sqrt{
\frac1{n_1}
+
\frac1{n_2}
}.
}
\]

%%%%%%%%%%%%%%%%%%%%%%%%%%%%%%%%%%%%%%%%%%%%%%%%%%%%%%%%%%%%
\subsection{Welch Versus Pooled Intervals}
\label{subsec:welch-versus-pooled}
%%%%%%%%%%%%%%%%%%%%%%%%%%%%%%%%%%%%%%%%%%%%%%%%%%%%%%%%%%%%

The pooled interval requires the additional assumption

\[
\boxed{
\sigma_1^2=\sigma_2^2.
}
\]

The Welch interval does not.

For this reason, Welch's method is often preferred as the default for
two independent means unless a common-variance model is scientifically
well justified.

%%%%%%%%%%%%%%%%%%%%%%%%%%%%%%%%%%%%%%%%%%%%%%%%%%%%%%%%%%%%
\subsection{Paired Confidence Interval}
\label{subsec:paired-confidence-interval}
%%%%%%%%%%%%%%%%%%%%%%%%%%%%%%%%%%%%%%%%%%%%%%%%%%%%%%%%%%%%

For paired observations, define

\[
D_i
=
X_i-Y_i.
\]

Then compute

\[
\overline{D}
\]

and

\[
S_D.
\]

The confidence interval for the population mean difference \(\mu_D\) is

\[
\boxed{
\overline{D}
\pm
t^*_{n-1}
\frac{
S_D
}{
\sqrt{n}
}.
}
\]

The paired problem is therefore treated as a one-sample problem on the
differences.

%%%%%%%%%%%%%%%%%%%%%%%%%%%%%%%%%%%%%%%%%%%%%%%%%%%%%%%%%%%%
\subsection{Worked Example: Paired Confidence Interval}
\label{subsec:worked-paired-ci}
%%%%%%%%%%%%%%%%%%%%%%%%%%%%%%%%%%%%%%%%%%%%%%%%%%%%%%%%%%%%

\begin{workedexamplebox}{Before-and-After Measurements}

Suppose \(n=16\) paired differences have

\[
\overline{d}=5
\]

and

\[
s_D=8.
\]

The standard error is

\[
\frac8{4}=2.
\]

For a \(95\%\) interval with \(15\) degrees of freedom, let the critical
value be \(t^*_{15}\).

Then

\begin{answerbox}
\[
\boxed{
5
\pm
2t^*_{15}.
}
\]
\end{answerbox}

\end{workedexamplebox}

%%%%%%%%%%%%%%%%%%%%%%%%%%%%%%%%%%%%%%%%%%%%%%%%%%%%%%%%%%%%
\subsection{Confidence Interval for Difference of Proportions}
\label{subsec:ci-difference-proportions}
%%%%%%%%%%%%%%%%%%%%%%%%%%%%%%%%%%%%%%%%%%%%%%%%%%%%%%%%%%%%

For independent sample proportions

\[
\widehat{p}_1
\]

and

\[
\widehat{p}_2,
\]

the parameter of interest is

\[
p_1-p_2.
\]

A large-sample confidence interval is

\[
\boxed{
(\widehat{p}_1-\widehat{p}_2)
\pm
z^*
\sqrt{
\frac{
\widehat{p}_1(1-\widehat{p}_1)
}{
n_1
}
+
\frac{
\widehat{p}_2(1-\widehat{p}_2)
}{
n_2
}
}.
}
\]

%%%%%%%%%%%%%%%%%%%%%%%%%%%%%%%%%%%%%%%%%%%%%%%%%%%%%%%%%%%%
\subsection{Worked Example: Difference of Proportions}
\label{subsec:worked-difference-proportions-ci}
%%%%%%%%%%%%%%%%%%%%%%%%%%%%%%%%%%%%%%%%%%%%%%%%%%%%%%%%%%%%

\begin{workedexamplebox}{Comparing Two Proportions}

Suppose

\[
\widehat{p}_1=0.70,
\qquad
n_1=200,
\]

and

\[
\widehat{p}_2=0.60,
\qquad
n_2=250.
\]

Then

\[
\widehat{p}_1-\widehat{p}_2
=
0.10.
\]

The estimated standard error is

\[
\sqrt{
\frac{
(0.70)(0.30)
}{200}
+
\frac{
(0.60)(0.40)
}{250}
}.
\]

Thus a \(95\%\) interval is

\begin{answerbox}
\[
\boxed{
0.10
\pm
1.96
\sqrt{
\frac{0.21}{200}
+
\frac{0.24}{250}
}.
}
\]
\end{answerbox}

\end{workedexamplebox}

%%%%%%%%%%%%%%%%%%%%%%%%%%%%%%%%%%%%%%%%%%%%%%%%%%%%%%%%%%%%
\subsection{Confidence Interval for a Normal Variance}
\label{subsec:ci-normal-variance}
%%%%%%%%%%%%%%%%%%%%%%%%%%%%%%%%%%%%%%%%%%%%%%%%%%%%%%%%%%%%

Suppose

\[
X_1,\ldots,X_n
\overset{\text{i.i.d.}}{\sim}
N(\mu,\sigma^2).
\]

Then

\[
\boxed{
\frac{
(n-1)S^2
}{
\sigma^2
}
\sim
\chi^2_{n-1}.
}
\]

Let

\[
\chi^2_{\alpha/2,\nu}
\]

and

\[
\chi^2_{1-\alpha/2,\nu}
\]

denote the lower and upper quantiles according to the convention being
used.

A two-sided interval for \(\sigma^2\) is

\[
\boxed{
\left(
\frac{
(n-1)S^2
}{
\chi^2_{1-\alpha/2,\nu}
},
\frac{
(n-1)S^2
}{
\chi^2_{\alpha/2,\nu}
}
\right),
}
\]

where

\[
\nu=n-1.
\]

%%%%%%%%%%%%%%%%%%%%%%%%%%%%%%%%%%%%%%%%%%%%%%%%%%%%%%%%%%%%
\subsection{Confidence Interval for a Standard Deviation}
\label{subsec:ci-normal-standard-deviation}
%%%%%%%%%%%%%%%%%%%%%%%%%%%%%%%%%%%%%%%%%%%%%%%%%%%%%%%%%%%%

After constructing an interval

\[
(L,U)
\]

for \(\sigma^2\), take square roots:

\[
\boxed{
(\sqrt{L},\sqrt{U})
}
\]

to obtain an interval for \(\sigma\).

Because the square-root function is increasing on positive values, the
coverage event is preserved.

%%%%%%%%%%%%%%%%%%%%%%%%%%%%%%%%%%%%%%%%%%%%%%%%%%%%%%%%%%%%
\subsection{Variance Interval Is Not Symmetric}
\label{subsec:variance-ci-asymmetry}
%%%%%%%%%%%%%%%%%%%%%%%%%%%%%%%%%%%%%%%%%%%%%%%%%%%%%%%%%%%%

The chi-square distribution is asymmetric.

Therefore confidence intervals for \(\sigma^2\) are generally not of the
simple form

\[
S^2\pm E.
\]

Instead, the lower and upper distances from \(S^2\) differ.

This is an important reminder that confidence intervals need not be
symmetric around the estimator.

%%%%%%%%%%%%%%%%%%%%%%%%%%%%%%%%%%%%%%%%%%%%%%%%%%%%%%%%%%%%
\subsection{Confidence Interval for Ratio of Normal Variances}
\label{subsec:ci-ratio-variances}
%%%%%%%%%%%%%%%%%%%%%%%%%%%%%%%%%%%%%%%%%%%%%%%%%%%%%%%%%%%%

Suppose two independent normal populations have variances

\[
\sigma_1^2
\]

and

\[
\sigma_2^2.
\]

Then

\[
\boxed{
\frac{
S_1^2/\sigma_1^2
}{
S_2^2/\sigma_2^2
}
\sim
F_{n_1-1,n_2-1}.
}
\]

This can be rearranged to produce a confidence interval for

\[
\boxed{
\frac{
\sigma_1^2
}{
\sigma_2^2
}.
}
\]

The exact endpoint formulas depend on the quantile notation adopted for
the \(F\)-distribution.

%%%%%%%%%%%%%%%%%%%%%%%%%%%%%%%%%%%%%%%%%%%%%%%%%%%%%%%%%%%%
\subsection{One-Sided Confidence Bounds}
\label{subsec:one-sided-confidence-bounds}
%%%%%%%%%%%%%%%%%%%%%%%%%%%%%%%%%%%%%%%%%%%%%%%%%%%%%%%%%%%%

Sometimes only an upper or lower bound is required.

A lower confidence bound has the form

\[
\boxed{
\theta\geq L.
}
\]

An upper confidence bound has the form

\[
\boxed{
\theta\leq U.
}
\]

For a normal mean with known \(\sigma\), a \(100(1-\alpha)\%\) lower
bound is

\[
\boxed{
\overline{X}
-
z_\alpha
\frac{\sigma}{\sqrt{n}}.
}
\]

Only one tail is allocated probability \(\alpha\).

%%%%%%%%%%%%%%%%%%%%%%%%%%%%%%%%%%%%%%%%%%%%%%%%%%%%%%%%%%%%
\subsection{Pivotal Quantities and Confidence Intervals}
\label{subsec:pivotal-confidence-intervals}
%%%%%%%%%%%%%%%%%%%%%%%%%%%%%%%%%%%%%%%%%%%%%%%%%%%%%%%%%%%%

A pivotal quantity has a distribution that does not depend on unknown
parameters.

Suppose

\[
Q(X,\theta)
\]

has known CDF independent of \(\theta\).

Find constants \(a\) and \(b\) such that

\[
P
(
a\leq Q(X,\theta)\leq b
)
=
1-\alpha.
\]

Then algebraically solve the inequalities for \(\theta\).

This is one of the most systematic methods for deriving exact confidence
intervals.

%%%%%%%%%%%%%%%%%%%%%%%%%%%%%%%%%%%%%%%%%%%%%%%%%%%%%%%%%%%%
\subsection{Worked Example: Pivot for a Normal Mean}
\label{subsec:worked-pivot-normal-mean}
%%%%%%%%%%%%%%%%%%%%%%%%%%%%%%%%%%%%%%%%%%%%%%%%%%%%%%%%%%%%

\begin{workedexamplebox}{Using a Standard Normal Pivot}

Suppose \(\sigma\) is known.

The pivot is

\[
Z
=
\frac{
\overline{X}-\mu
}{
\sigma/\sqrt{n}
}.
\]

Choose \(z^*\) so that

\[
P(-z^*\leq Z\leq z^*)=1-\alpha.
\]

Then

\[
-z^*
\leq
\frac{
\overline{X}-\mu
}{
\sigma/\sqrt{n}
}
\leq
z^*.
\]

Solving for \(\mu\),

\begin{answerbox}
\[
\boxed{
\mu
\in
\left[
\overline{X}
-
z^*\frac{\sigma}{\sqrt{n}},
\,
\overline{X}
+
z^*\frac{\sigma}{\sqrt{n}}
\right].
}
\]
\end{answerbox}

\end{workedexamplebox}

%%%%%%%%%%%%%%%%%%%%%%%%%%%%%%%%%%%%%%%%%%%%%%%%%%%%%%%%%%%%
\subsection{Wald Confidence Intervals}
\label{subsec:wald-confidence-intervals}
%%%%%%%%%%%%%%%%%%%%%%%%%%%%%%%%%%%%%%%%%%%%%%%%%%%%%%%%%%%%

Suppose an estimator is approximately normal:

\[
\widehat{\theta}
\approx
N
(
\theta,
\operatorname{SE}^2
).
\]

Then the generic Wald interval is

\[
\boxed{
\widehat{\theta}
\pm
z^*
\widehat{\operatorname{SE}}
(
\widehat{\theta}
).
}
\]

Wald intervals are simple and widespread.

However, they can perform poorly when:

\[
\boxed{
\text{sample sizes are small},
}
\]

\[
\boxed{
\text{parameters are near boundaries},
}
\]

or

\[
\boxed{
\text{sampling distributions are strongly skewed}.
}
\]

%%%%%%%%%%%%%%%%%%%%%%%%%%%%%%%%%%%%%%%%%%%%%%%%%%%%%%%%%%%%
\subsection{Likelihood-Ratio Confidence Intervals}
\label{subsec:likelihood-ratio-ci}
%%%%%%%%%%%%%%%%%%%%%%%%%%%%%%%%%%%%%%%%%%%%%%%%%%%%%%%%%%%%

Let

\[
\widehat{\theta}
\]

be the MLE.

A likelihood-ratio confidence set can be based on parameter values for
which

\[
\boxed{
2
\left[
\ell(\widehat{\theta})
-
\ell(\theta)
\right]
}
\]

is not too large.

Under regularity conditions,

\[
\boxed{
2
\left[
\ell(\widehat{\theta})
-
\ell(\theta)
\right]
\approx
\chi^2_1
}
\]

for a scalar parameter.

Thus an approximate \(1-\alpha\) confidence set is

\[
\boxed{
\left\{
\theta:
2
[
\ell(\widehat{\theta})
-
\ell(\theta)
]
\leq
\chi^2_{1-\alpha,1}
\right\}.
}
\]

%%%%%%%%%%%%%%%%%%%%%%%%%%%%%%%%%%%%%%%%%%%%%%%%%%%%%%%%%%%%
\subsection{Why Likelihood Intervals Can Be Asymmetric}
\label{subsec:likelihood-ci-asymmetry}
%%%%%%%%%%%%%%%%%%%%%%%%%%%%%%%%%%%%%%%%%%%%%%%%%%%%%%%%%%%%

Likelihood functions are not always symmetric around their maximum.

Therefore a likelihood-based interval may naturally extend farther in
one direction than the other.

This can be useful for:

\[
\boxed{
\text{skewed parameters},
}
\]

\[
\boxed{
\text{positive parameters},
}
\]

or

\[
\boxed{
\text{parameters near boundaries}.
}
\]

%%%%%%%%%%%%%%%%%%%%%%%%%%%%%%%%%%%%%%%%%%%%%%%%%%%%%%%%%%%%
\subsection{Score Confidence Intervals}
\label{subsec:score-confidence-intervals}
%%%%%%%%%%%%%%%%%%%%%%%%%%%%%%%%%%%%%%%%%%%%%%%%%%%%%%%%%%%%

Score intervals are obtained by inverting a score test.

The Wilson interval for a binomial proportion is an important example.

Score intervals often improve on simple Wald intervals because the
standardization is performed under each candidate parameter value rather
than only at the observed estimate.

%%%%%%%%%%%%%%%%%%%%%%%%%%%%%%%%%%%%%%%%%%%%%%%%%%%%%%%%%%%%
\subsection{Transformation-Based Confidence Intervals}
\label{subsec:transformation-confidence-intervals}
%%%%%%%%%%%%%%%%%%%%%%%%%%%%%%%%%%%%%%%%%%%%%%%%%%%%%%%%%%%%

Suppose

\[
\eta=g(\theta)
\]

and \(g\) is monotone.

If

\[
(L,U)
\]

is a confidence interval for \(\theta\), then

\[
\boxed{
(g(L),g(U))
}
\]

is a confidence interval for \(\eta\) when \(g\) is increasing.

If \(g\) is decreasing, the endpoint order reverses.

This provides a simple way to obtain intervals for transformed
parameters.

%%%%%%%%%%%%%%%%%%%%%%%%%%%%%%%%%%%%%%%%%%%%%%%%%%%%%%%%%%%%
\subsection{Log-Scale Confidence Intervals}
\label{subsec:log-scale-confidence-intervals}
%%%%%%%%%%%%%%%%%%%%%%%%%%%%%%%%%%%%%%%%%%%%%%%%%%%%%%%%%%%%

Positive parameters are often estimated on the logarithmic scale.

Suppose

\[
\log\theta
\]

has approximate interval

\[
(L,U).
\]

Exponentiating gives

\[
\boxed{
(e^L,e^U)
}
\]

for \(\theta\).

The resulting interval is automatically positive and typically
asymmetric on the original scale.

%%%%%%%%%%%%%%%%%%%%%%%%%%%%%%%%%%%%%%%%%%%%%%%%%%%%%%%%%%%%
\subsection{Delta-Method Confidence Intervals}
\label{subsec:delta-method-confidence-intervals}
%%%%%%%%%%%%%%%%%%%%%%%%%%%%%%%%%%%%%%%%%%%%%%%%%%%%%%%%%%%%

Suppose

\[
\widehat{\theta}
\]

is approximately normal with variance \(V/n\).

For differentiable \(g\),

\[
g(\widehat{\theta})
\]

has approximate variance

\[
\boxed{
[g'(\theta)]^2
\frac{V}{n}.
}
\]

Replacing unknown quantities by estimates gives a Wald-type interval for

\[
g(\theta).
\]

This is useful when the parameter of interest is a transformation of an
easier estimator.

%%%%%%%%%%%%%%%%%%%%%%%%%%%%%%%%%%%%%%%%%%%%%%%%%%%%%%%%%%%%
\subsection{Bootstrap Confidence Intervals}
\label{subsec:bootstrap-confidence-intervals}
%%%%%%%%%%%%%%%%%%%%%%%%%%%%%%%%%%%%%%%%%%%%%%%%%%%%%%%%%%%%

Suppose bootstrap replicates are

\[
\widehat{\theta}_1^*,
\ldots,
\widehat{\theta}_B^*.
\]

Their empirical distribution approximates the sampling distribution of
the estimator.

This can be used to construct confidence intervals when an analytic
standard error or exact sampling distribution is difficult to obtain.

%%%%%%%%%%%%%%%%%%%%%%%%%%%%%%%%%%%%%%%%%%%%%%%%%%%%%%%%%%%%
\subsection{Bootstrap Normal Interval}
\label{subsec:bootstrap-normal-interval}
%%%%%%%%%%%%%%%%%%%%%%%%%%%%%%%%%%%%%%%%%%%%%%%%%%%%%%%%%%%%

If the bootstrap distribution is approximately symmetric and normal,
one may use

\[
\boxed{
\widehat{\theta}
\pm
z^*
\widehat{\operatorname{SE}}_{\text{boot}}.
}
\]

This mirrors the ordinary Wald interval but uses a bootstrap estimate of
standard error.

%%%%%%%%%%%%%%%%%%%%%%%%%%%%%%%%%%%%%%%%%%%%%%%%%%%%%%%%%%%%
\subsection{Bootstrap Percentile Interval}
\label{subsec:bootstrap-percentile-interval}
%%%%%%%%%%%%%%%%%%%%%%%%%%%%%%%%%%%%%%%%%%%%%%%%%%%%%%%%%%%%

Let

\[
q_{\alpha/2}^*
\]

and

\[
q_{1-\alpha/2}^*
\]

be empirical quantiles of the bootstrap estimates.

The percentile interval is

\[
\boxed{
\left(
q_{\alpha/2}^*,
q_{1-\alpha/2}^*
\right).
}
\]

It is simple and automatically reflects some asymmetry in the bootstrap
distribution.

%%%%%%%%%%%%%%%%%%%%%%%%%%%%%%%%%%%%%%%%%%%%%%%%%%%%%%%%%%%%
\subsection{Basic Bootstrap Interval}
\label{subsec:basic-bootstrap-interval}
%%%%%%%%%%%%%%%%%%%%%%%%%%%%%%%%%%%%%%%%%%%%%%%%%%%%%%%%%%%%

The basic bootstrap interval reflects the percentile deviations around
the observed estimate.

If

\[
q_{\alpha/2}^*
\]

and

\[
q_{1-\alpha/2}^*
\]

are bootstrap quantiles, the basic interval is

\[
\boxed{
\left(
2\widehat{\theta}
-
q_{1-\alpha/2}^*,
\,
2\widehat{\theta}
-
q_{\alpha/2}^*
\right).
}
\]

%%%%%%%%%%%%%%%%%%%%%%%%%%%%%%%%%%%%%%%%%%%%%%%%%%%%%%%%%%%%
\subsection{Bias-Corrected Bootstrap Preview}
\label{subsec:bca-bootstrap-preview}
%%%%%%%%%%%%%%%%%%%%%%%%%%%%%%%%%%%%%%%%%%%%%%%%%%%%%%%%%%%%

More advanced bootstrap procedures adjust for:

\[
\boxed{
\text{bias},
}
\]

and

\[
\boxed{
\text{skewness}.
}
\]

The bias-corrected and accelerated interval, abbreviated BCa, is a widely
used higher-order bootstrap interval.

Its full derivation belongs to computational statistics.

%%%%%%%%%%%%%%%%%%%%%%%%%%%%%%%%%%%%%%%%%%%%%%%%%%%%%%%%%%%%
\subsection{Bootstrap Limitations for Confidence Intervals}
\label{subsec:bootstrap-ci-limitations}
%%%%%%%%%%%%%%%%%%%%%%%%%%%%%%%%%%%%%%%%%%%%%%%%%%%%%%%%%%%%

Bootstrap confidence intervals may fail or perform poorly for:

\[
\boxed{
\text{very small samples},
}
\]

\[
\boxed{
\text{extreme-value estimators},
}
\]

\[
\boxed{
\text{strong dependence without appropriate resampling},
}
\]

\[
\boxed{
\text{nonsmooth statistics},
}
\]

or

\[
\boxed{
\text{poor sampling designs}.
}
\]

Resampling does not remove systematic bias in data collection.

%%%%%%%%%%%%%%%%%%%%%%%%%%%%%%%%%%%%%%%%%%%%%%%%%%%%%%%%%%%%
\subsection{Confidence Intervals Under Dependence}
\label{subsec:ci-under-dependence}
%%%%%%%%%%%%%%%%%%%%%%%%%%%%%%%%%%%%%%%%%%%%%%%%%%%%%%%%%%%%

For dependent observations, ordinary independent-sample standard errors
may be incorrect.

Examples include:

\[
\boxed{
\text{time-series data},
}
\]

\[
\boxed{
\text{clustered observations},
}
\]

and

\[
\boxed{
\text{repeated measurements}.
}
\]

The confidence interval must use a variance estimator appropriate to the
dependence structure.

%%%%%%%%%%%%%%%%%%%%%%%%%%%%%%%%%%%%%%%%%%%%%%%%%%%%%%%%%%%%
\subsection{Cluster-Robust Confidence Intervals Preview}
\label{subsec:cluster-robust-ci-preview}
%%%%%%%%%%%%%%%%%%%%%%%%%%%%%%%%%%%%%%%%%%%%%%%%%%%%%%%%%%%%

When observations are independent across clusters but dependent within
clusters, one may estimate covariance using cluster-robust methods.

The generic interval remains

\[
\boxed{
\text{estimate}
\pm
\text{critical value}
\times
\text{cluster-adjusted SE}.
}
\]

The key change is the standard-error estimator.

%%%%%%%%%%%%%%%%%%%%%%%%%%%%%%%%%%%%%%%%%%%%%%%%%%%%%%%%%%%%
\subsection{Confidence Intervals from Complex Surveys}
\label{subsec:complex-survey-confidence-intervals}
%%%%%%%%%%%%%%%%%%%%%%%%%%%%%%%%%%%%%%%%%%%%%%%%%%%%%%%%%%%%

Survey estimates may involve:

\[
\boxed{
\text{weights},
}
\]

\[
\boxed{
\text{strata},
}
\]

and

\[
\boxed{
\text{clusters}.
}
\]

Standard errors must reflect the sampling design.

Applying simple-random-sample formulas to complex survey data may give
incorrect uncertainty estimates.

%%%%%%%%%%%%%%%%%%%%%%%%%%%%%%%%%%%%%%%%%%%%%%%%%%%%%%%%%%%%
\subsection{Confidence Interval Width}
\label{subsec:confidence-interval-width}
%%%%%%%%%%%%%%%%%%%%%%%%%%%%%%%%%%%%%%%%%%%%%%%%%%%%%%%%%%%%

For a symmetric interval

\[
\widehat{\theta}
\pm E,
\]

the total width is

\[
\boxed{
2E.
}
\]

For a normal mean with known variance,

\[
\boxed{
\text{width}
=
2z^*
\frac{\sigma}{\sqrt{n}}.
}
\]

Thus width decreases at rate

\[
\boxed{
n^{-1/2}.
}
\]

%%%%%%%%%%%%%%%%%%%%%%%%%%%%%%%%%%%%%%%%%%%%%%%%%%%%%%%%%%%%
\subsection{Doubling Precision Is Expensive}
\label{subsec:precision-sample-size}
%%%%%%%%%%%%%%%%%%%%%%%%%%%%%%%%%%%%%%%%%%%%%%%%%%%%%%%%%%%%

Suppose the margin of error is proportional to

\[
\frac1{\sqrt{n}}.
\]

To cut the margin of error in half, solve

\[
\frac1{\sqrt{n_{\text{new}}}}
=
\frac12
\frac1{\sqrt{n_{\text{old}}}}.
\]

Then

\[
\boxed{
n_{\text{new}}
=
4n_{\text{old}}.
}
\]

Thus doubling statistical precision often requires approximately four
times as much data.

%%%%%%%%%%%%%%%%%%%%%%%%%%%%%%%%%%%%%%%%%%%%%%%%%%%%%%%%%%%%
\subsection{Confidence Interval Versus Prediction Interval}
\label{subsec:confidence-versus-prediction-interval}
%%%%%%%%%%%%%%%%%%%%%%%%%%%%%%%%%%%%%%%%%%%%%%%%%%%%%%%%%%%%

A confidence interval estimates an unknown population parameter such as

\[
\mu.
\]

A prediction interval aims to contain a future observation.

Prediction intervals are typically wider because they must account for:

\[
\boxed{
\text{uncertainty in the estimated mean}
}
\]

and

\[
\boxed{
\text{individual observation variability}.
}
\]

Prediction intervals are developed further in regression.

%%%%%%%%%%%%%%%%%%%%%%%%%%%%%%%%%%%%%%%%%%%%%%%%%%%%%%%%%%%%
\subsection{Confidence Interval Versus Tolerance Interval}
\label{subsec:confidence-versus-tolerance}
%%%%%%%%%%%%%%%%%%%%%%%%%%%%%%%%%%%%%%%%%%%%%%%%%%%%%%%%%%%%

A tolerance interval aims to contain a specified proportion of a
population with a stated confidence level.

For example, one might seek an interval that contains at least

\[
95\%
\]

of the population with

\[
99\%
\]

confidence.

This is different from a confidence interval for a parameter.

%%%%%%%%%%%%%%%%%%%%%%%%%%%%%%%%%%%%%%%%%%%%%%%%%%%%%%%%%%%%
\subsection{Confidence Intervals and Hypothesis Tests}
\label{subsec:confidence-test-duality}
%%%%%%%%%%%%%%%%%%%%%%%%%%%%%%%%%%%%%%%%%%%%%%%%%%%%%%%%%%%%

A two-sided confidence interval and a two-sided hypothesis test are often
dual.

Suppose

\[
H_0:\theta=\theta_0
\]

is tested at significance level \(\alpha\).

Under many standard procedures,

\[
\boxed{
\theta_0
\text{ is rejected}
}
\]

if and only if

\[
\boxed{
\theta_0
\text{ lies outside the }100(1-\alpha)\%\text{ confidence interval}.
}
\]

This duality becomes central in the next section.

%%%%%%%%%%%%%%%%%%%%%%%%%%%%%%%%%%%%%%%%%%%%%%%%%%%%%%%%%%%%
\subsection{Confidence Intervals and Practical Significance}
\label{subsec:ci-practical-significance}
%%%%%%%%%%%%%%%%%%%%%%%%%%%%%%%%%%%%%%%%%%%%%%%%%%%%%%%%%%%%

A confidence interval conveys more information than a binary
significance decision.

Suppose an estimated treatment effect is

\[
2.0
\]

with interval

\[
(1.8,2.2).
\]

The effect is estimated precisely.

But whether an effect near \(2\) is practically important depends on
scientific context.

Thus confidence intervals should be interpreted in terms of:

\[
\boxed{
\text{magnitude}
+
\text{precision}
+
\text{scientific relevance}.
}
\]

%%%%%%%%%%%%%%%%%%%%%%%%%%%%%%%%%%%%%%%%%%%%%%%%%%%%%%%%%%%%
\subsection{Narrow Intervals Can Still Be Biased}
\label{subsec:narrow-ci-biased}
%%%%%%%%%%%%%%%%%%%%%%%%%%%%%%%%%%%%%%%%%%%%%%%%%%%%%%%%%%%%

A very large sample can produce a narrow interval around a biased
estimate.

For example, severe undercoverage in a survey may produce a precise
estimate of the wrong population quantity.

Confidence intervals generally quantify sampling uncertainty under the
assumed design and model.

They do not automatically account for every form of:

\[
\boxed{
\text{selection bias},
}
\]

\[
\boxed{
\text{measurement bias},
}
\]

or

\[
\boxed{
\text{model misspecification}.
}
\]

%%%%%%%%%%%%%%%%%%%%%%%%%%%%%%%%%%%%%%%%%%%%%%%%%%%%%%%%%%%%
\subsection{Confidence Level Does Not Measure Model Validity}
\label{subsec:confidence-not-model-validity}
%%%%%%%%%%%%%%%%%%%%%%%%%%%%%%%%%%%%%%%%%%%%%%%%%%%%%%%%%%%%

A \(95\%\) confidence interval does not mean there is a \(95\%\) chance
that all modeling assumptions are correct.

Coverage statements are conditional on the procedure and its
assumptions.

If assumptions fail badly, actual coverage can differ from the nominal
confidence level.

%%%%%%%%%%%%%%%%%%%%%%%%%%%%%%%%%%%%%%%%%%%%%%%%%%%%%%%%%%%%
\subsection{Multiple Confidence Intervals}
\label{subsec:multiple-confidence-intervals-preview}
%%%%%%%%%%%%%%%%%%%%%%%%%%%%%%%%%%%%%%%%%%%%%%%%%%%%%%%%%%%%

Suppose \(m\) separate \(95\%\) intervals are constructed.

The probability that all intervals simultaneously contain their
parameters is generally less than \(95\%\).

This motivates simultaneous confidence procedures.

A simple conservative method is Bonferroni adjustment.

%%%%%%%%%%%%%%%%%%%%%%%%%%%%%%%%%%%%%%%%%%%%%%%%%%%%%%%%%%%%
\subsection{Bonferroni Confidence Intervals Preview}
\label{subsec:bonferroni-ci-preview}
%%%%%%%%%%%%%%%%%%%%%%%%%%%%%%%%%%%%%%%%%%%%%%%%%%%%%%%%%%%%

To obtain overall confidence approximately at least \(1-\alpha\) for
\(m\) intervals, construct each individual interval at level

\[
\boxed{
1-\frac{\alpha}{m}.
}
\]

By the union bound, the probability that any interval misses its target
is at most

\[
m\frac{\alpha}{m}
=
\alpha.
\]

Thus simultaneous coverage is at least \(1-\alpha\).

%%%%%%%%%%%%%%%%%%%%%%%%%%%%%%%%%%%%%%%%%%%%%%%%%%%%%%%%%%%%
\subsection{Confidence Regions}
\label{subsec:confidence-regions}
%%%%%%%%%%%%%%%%%%%%%%%%%%%%%%%%%%%%%%%%%%%%%%%%%%%%%%%%%%%%

For a parameter vector

\[
\boldsymbol{\theta}
\in
\mathbb{R}^k,
\]

a confidence set need not be an interval.

It may be a region such as

\[
\boxed{
\mathcal{C}(X)
\subseteq
\mathbb{R}^k.
}
\]

For approximately multivariate normal estimators, confidence regions are
often ellipsoidal.

%%%%%%%%%%%%%%%%%%%%%%%%%%%%%%%%%%%%%%%%%%%%%%%%%%%%%%%%%%%%
\subsection{Ellipsoidal Confidence Region Preview}
\label{subsec:ellipsoidal-confidence-region}
%%%%%%%%%%%%%%%%%%%%%%%%%%%%%%%%%%%%%%%%%%%%%%%%%%%%%%%%%%%%

Suppose

\[
\widehat{\boldsymbol{\theta}}
\]

is approximately multivariate normal with covariance matrix

\[
\Sigma.
\]

Then a quadratic form

\[
\boxed{
(
\widehat{\boldsymbol{\theta}}
-
\boldsymbol{\theta}
)^T
\Sigma^{-1}
(
\widehat{\boldsymbol{\theta}}
-
\boldsymbol{\theta}
)
}
\]

is approximately chi-square distributed.

This leads to confidence ellipsoids in multiparameter problems.

%%%%%%%%%%%%%%%%%%%%%%%%%%%%%%%%%%%%%%%%%%%%%%%%%%%%%%%%%%%%
\subsection{Confidence Interval Workflow}
\label{subsec:confidence-interval-workflow}
%%%%%%%%%%%%%%%%%%%%%%%%%%%%%%%%%%%%%%%%%%%%%%%%%%%%%%%%%%%%

A systematic confidence-interval workflow is:

\begin{enumerate}
    \item identify the parameter of interest;
    \item identify an estimator;
    \item identify the sampling design or probability model;
    \item determine the estimator's standard error;
    \item determine whether an exact pivot is available;
    \item otherwise choose an appropriate asymptotic approximation;
    \item select the confidence level;
    \item obtain the critical value;
    \item calculate the interval;
    \item check assumptions;
    \item interpret the interval in context;
    \item report the estimate, interval, confidence level, and relevant
          sample size.
\end{enumerate}

%%%%%%%%%%%%%%%%%%%%%%%%%%%%%%%%%%%%%%%%%%%%%%%%%%%%%%%%%%%%
\subsection{Choosing the Correct Interval}
\label{subsec:choosing-correct-ci}
%%%%%%%%%%%%%%%%%%%%%%%%%%%%%%%%%%%%%%%%%%%%%%%%%%%%%%%%%%%%

For one quantitative population mean:

\[
\boxed{
\text{known }\sigma
\rightarrow
z\text{-interval},
}
\]

\[
\boxed{
\text{unknown }\sigma
\rightarrow
t\text{-interval}.
}
\]

For one proportion:

\[
\boxed{
\text{binomial proportion interval}.
}
\]

For two independent quantitative groups:

\[
\boxed{
\text{Welch or pooled two-sample interval}.
}
\]

For paired data:

\[
\boxed{
\text{one-sample interval on differences}.
}
\]

For normal variance:

\[
\boxed{
\chi^2\text{-interval}.
}
\]

%%%%%%%%%%%%%%%%%%%%%%%%%%%%%%%%%%%%%%%%%%%%%%%%%%%%%%%%%%%%
\subsection{Common Errors}
\label{subsec:confidence-interval-common-errors}
%%%%%%%%%%%%%%%%%%%%%%%%%%%%%%%%%%%%%%%%%%%%%%%%%%%%%%%%%%%%

\subsubsection{Assigning Frequentist Probability to a Fixed Parameter}

Once the interval is observed, the frequentist parameter is fixed.

The confidence level describes the long-run success rate of the
procedure.

%%%%%%%%%%%%%%%%%%%%%%%%%%%%%%%%%%%%%%%%%%%%%%%%%%%%%%%%%%%%
\subsubsection{Confusing Confidence Level with Percentage of Data}

A \(95\%\) confidence interval for \(\mu\) does not mean the interval
contains \(95\%\) of individual observations.

It is an interval for a parameter.

%%%%%%%%%%%%%%%%%%%%%%%%%%%%%%%%%%%%%%%%%%%%%%%%%%%%%%%%%%%%
\subsubsection{Confusing Confidence and Prediction Intervals}

A confidence interval estimates a population feature.

A prediction interval targets a future observation.

%%%%%%%%%%%%%%%%%%%%%%%%%%%%%%%%%%%%%%%%%%%%%%%%%%%%%%%%%%%%
\subsubsection{Using \(z\) Instead of \(t\) with Unknown \(\sigma\)}

For exact inference on a normal mean when \(\sigma\) is unknown, use

\[
\boxed{
t_{n-1}.
}
\]

%%%%%%%%%%%%%%%%%%%%%%%%%%%%%%%%%%%%%%%%%%%%%%%%%%%%%%%%%%%%
\subsubsection{Using \(t\) Because the Sample Is Small Without Considering the Model}

The exact \(t\)-distribution comes from normal sampling with estimated
variance.

It is not simply a generic ``small-sample distribution.''

%%%%%%%%%%%%%%%%%%%%%%%%%%%%%%%%%%%%%%%%%%%%%%%%%%%%%%%%%%%%
\subsubsection{Ignoring Pairing}

For matched or repeated observations, use differences.

Treating paired data as independent discards within-pair information.

%%%%%%%%%%%%%%%%%%%%%%%%%%%%%%%%%%%%%%%%%%%%%%%%%%%%%%%%%%%%
\subsubsection{Pooling Variances Automatically}

The pooled two-sample interval assumes equal population variances.

Welch's method does not require that equality assumption.

%%%%%%%%%%%%%%%%%%%%%%%%%%%%%%%%%%%%%%%%%%%%%%%%%%%%%%%%%%%%
\subsubsection{Using the Wald Proportion Interval Near Zero or One}

The simple Wald interval can have poor coverage and may even extend
outside

\[
[0,1].
\]

Wilson or exact methods are often preferable.

%%%%%%%%%%%%%%%%%%%%%%%%%%%%%%%%%%%%%%%%%%%%%%%%%%%%%%%%%%%%
\subsubsection{Reporting More Decimal Places Than the Data Support}

A confidence interval should not imply more numerical precision than the
measurement process justifies.

%%%%%%%%%%%%%%%%%%%%%%%%%%%%%%%%%%%%%%%%%%%%%%%%%%%%%%%%%%%%
\subsubsection{Ignoring Study Design}

A mathematically correct interval formula applied to a biased or
incorrect sampling design may still produce misleading conclusions.

%%%%%%%%%%%%%%%%%%%%%%%%%%%%%%%%%%%%%%%%%%%%%%%%%%%%%%%%%%%%
\subsubsection{Interpreting Nonoverlapping Intervals as the Only Way to Compare Groups}

Two separate confidence intervals are not always the correct tool for
testing or estimating a difference.

Construct the interval directly for

\[
\mu_1-\mu_2
\]

or

\[
p_1-p_2.
\]

%%%%%%%%%%%%%%%%%%%%%%%%%%%%%%%%%%%%%%%%%%%%%%%%%%%%%%%%%%%%
\subsubsection{Assuming Overlapping \(95\%\) Intervals Means No Significant Difference}

Two individual \(95\%\) confidence intervals can overlap even when a
formal comparison of the difference is statistically significant.

The direct sampling distribution of the difference should be used.

%%%%%%%%%%%%%%%%%%%%%%%%%%%%%%%%%%%%%%%%%%%%%%%%%%%%%%%%%%%%
\subsubsection{Treating Confidence Width as Only a Function of Sample Size}

Width also depends on:

\[
\boxed{
\text{confidence level},
}
\]

\[
\boxed{
\text{variability},
}
\]

and

\[
\boxed{
\text{sampling design}.
}
\]

%%%%%%%%%%%%%%%%%%%%%%%%%%%%%%%%%%%%%%%%%%%%%%%%%%%%%%%%%%%%
\subsection{Applications}
\label{subsec:confidence-interval-applications}
%%%%%%%%%%%%%%%%%%%%%%%%%%%%%%%%%%%%%%%%%%%%%%%%%%%%%%%%%%%%

\begin{applicationbox}

\textbf{Survey Research.}
Confidence intervals quantify sampling uncertainty in estimated means,
proportions, and population differences.

\medskip

\textbf{Scientific Experiments.}
Intervals describe plausible treatment-effect magnitudes rather than
only whether an effect is statistically detectable.

\medskip

\textbf{Engineering.}
Intervals estimate process means, defect rates, lifetimes, and
variability.

\medskip

\textbf{Environmental Science.}
Researchers use intervals for pollutant concentrations, temperature
changes, ecological measurements, and model parameters.

\medskip

\textbf{Medicine.}
Treatment effects, response proportions, risk differences, and mean
differences are commonly reported with confidence intervals.

\medskip

\textbf{Business and Economics.}
Confidence intervals quantify uncertainty in estimated demand,
differences, regression coefficients, and economic indicators.

\medskip

\textbf{Data Science.}
Asymptotic, bootstrap, and model-based intervals quantify uncertainty in
estimated model parameters and performance metrics.

\medskip

\textbf{Public Reporting.}
Intervals communicate statistical precision more effectively than point
estimates alone.

\end{applicationbox}

%%%%%%%%%%%%%%%%%%%%%%%%%%%%%%%%%%%%%%%%%%%%%%%%%%%%%%%%%%%%
\subsection{Exercises}
\label{subsec:confidence-interval-exercises}
%%%%%%%%%%%%%%%%%%%%%%%%%%%%%%%%%%%%%%%%%%%%%%%%%%%%%%%%%%%%

\begin{exercise}[1]
Define a confidence interval.
\end{exercise}

\begin{exercise}[1]
What does a \(95\%\) confidence level mean in repeated sampling?
\end{exercise}

\begin{exercise}[1]
Define margin of error.
\end{exercise}

\begin{exercise}[1]
State the known-\(\sigma\) confidence interval for a population mean.
\end{exercise}

\begin{exercise}[1]
State the one-sample \(t\)-interval.
\end{exercise}

\begin{exercise}[1]
State the large-sample Wald confidence interval for a proportion.
\end{exercise}

\begin{exercise}[1]
State the Welch confidence interval for a difference of means.
\end{exercise}

\begin{exercise}[1]
State the paired confidence interval.
\end{exercise}

\begin{exercise}[1]
State the large-sample interval for a difference of proportions.
\end{exercise}

\begin{exercise}[1]
State the chi-square pivot used for a normal variance interval.
\end{exercise}

\begin{exercise}[2]
Suppose

\[
\overline{x}=40,
\qquad
\sigma=8,
\qquad
n=64.
\]

Construct a \(95\%\) confidence interval for \(\mu\).
\end{exercise}

\begin{exercise}[2]
For the same data, construct a \(99\%\) confidence interval and compare
its width with the \(95\%\) interval.
\end{exercise}

\begin{exercise}[2]
Suppose

\[
n=36,
\qquad
\overline{x}=25,
\qquad
s=6.
\]

Write the \(95\%\) one-sample \(t\)-interval in terms of the appropriate
critical value.
\end{exercise}

\begin{exercise}[2]
Suppose

\[
n=500,
\qquad
x=300.
\]

Construct the large-sample Wald \(95\%\) confidence interval for \(p\).
\end{exercise}

\begin{exercise}[2]
For the previous data, write the Wilson interval formula and compare its
center with \(\widehat{p}\).
\end{exercise}

\begin{exercise}[2]
Find the required sample size to estimate a mean with margin of error at
most \(3\), assuming

\[
\sigma=15
\]

and \(95\%\) confidence.
\end{exercise}

\begin{exercise}[2]
Find the conservative sample size required to estimate a population
proportion with margin of error \(0.04\) at \(95\%\) confidence.
\end{exercise}

\begin{exercise}[2]
Suppose

\[
\overline{x}_1=20,
\quad
s_1=5,
\quad
n_1=30,
\]

and

\[
\overline{x}_2=17,
\quad
s_2=4,
\quad
n_2=35.
\]

Write the Welch confidence interval for \(\mu_1-\mu_2\).
\end{exercise}

\begin{exercise}[2]
Suppose paired differences have

\[
\overline{d}=2.5,
\qquad
s_D=4,
\qquad
n=25.
\]

Construct the \(95\%\) paired interval in terms of the appropriate
\(t\)-critical value.
\end{exercise}

\begin{exercise}[2]
Suppose

\[
\widehat{p}_1=0.55,
\quad
n_1=200,
\]

and

\[
\widehat{p}_2=0.45,
\quad
n_2=250.
\]

Construct a \(95\%\) interval for \(p_1-p_2\).
\end{exercise}

\begin{exercise}[3]
Derive the known-\(\sigma\) interval for a normal mean by inverting the
standard normal pivot.
\end{exercise}

\begin{exercise}[3]
Derive the one-sample \(t\)-interval from the \(t\)-pivot.
\end{exercise}

\begin{exercise}[3]
Show that increasing \(n\) by a factor of \(4\) halves the known-\(\sigma\)
margin of error.
\end{exercise}

\begin{exercise}[3]
Derive the sample-size formula

\[
n
=
\left(
\frac{
z^*\sigma
}{E}
\right)^2.
\]
\end{exercise}

\begin{exercise}[3]
Show that

\[
p(1-p)
\]

is maximized at

\[
p=\frac12.
\]
\end{exercise}

\begin{exercise}[3]
Derive the conservative sample-size formula for a proportion by using

\[
p=0.5.
\]
\end{exercise}

\begin{exercise}[3]
Explain why the Welch interval is often preferred to the pooled interval.
\end{exercise}

\begin{exercise}[3]
Derive the pooled variance estimator from the two independent sample sums
of squares.
\end{exercise}

\begin{exercise}[3]
Explain why a paired interval uses the standard deviation of the
differences rather than separate standard deviations.
\end{exercise}

\begin{exercise}[3]
Derive the large-sample standard error for the difference of two
independent sample proportions.
\end{exercise}

\begin{exercise}[3]
Derive a confidence interval for a normal variance from the chi-square
pivot.
\end{exercise}

\begin{exercise}[3]
Explain why the variance interval is asymmetric.
\end{exercise}

\begin{exercise}[3]
Derive a confidence interval for the ratio of two normal variances using
the \(F\)-distribution.
\end{exercise}

\begin{exercise}[3]
Show how a monotone transformation converts a confidence interval for
\(\theta\) into an interval for \(g(\theta)\).
\end{exercise}

\begin{exercise}[3]
Use the delta method to construct an approximate confidence interval for

\[
\log\theta.
\]
\end{exercise}

\begin{exercise}[3]
Explain why likelihood-ratio intervals may be asymmetric.
\end{exercise}

\begin{exercise}[3]
Explain the difference between a Wald interval and a score interval.
\end{exercise}

\begin{exercise}[3]
Construct a basic bootstrap interval from given bootstrap quantiles.
\end{exercise}

\begin{exercise}[3]
Compare percentile and basic bootstrap confidence intervals.
\end{exercise}

\begin{exercise}[3]
Explain why bootstrapping cannot repair a nonrepresentative sample.
\end{exercise}

\begin{exercise}[3]
Explain the duality between a two-sided \(1-\alpha\) confidence interval
and a level-\(\alpha\) hypothesis test.
\end{exercise}

\begin{exercise}[3]
Explain why comparing two separate group confidence intervals is not the
same as constructing an interval for their difference.
\end{exercise}

\begin{exercise}[3]
Explain why a narrow confidence interval does not guarantee an unbiased
study.
\end{exercise}

\begin{exercise}[4]
Study exact confidence intervals derived from pivotal quantities.
\end{exercise}

\begin{exercise}[4]
Investigate the coverage probability of the Wald binomial interval.
\end{exercise}

\begin{exercise}[4]
Derive the Wilson score interval algebraically.
\end{exercise}

\begin{exercise}[4]
Study the Clopper--Pearson exact binomial interval.
\end{exercise}

\begin{exercise}[4]
Compare Wald, Wilson, Agresti--Coull, and exact binomial intervals through
simulation.
\end{exercise}

\begin{exercise}[4]
Study the Welch--Satterthwaite approximation.
\end{exercise}

\begin{exercise}[4]
Investigate likelihood-ratio confidence intervals and Wilks' theorem.
\end{exercise}

\begin{exercise}[4]
Study score confidence intervals in regular parametric models.
\end{exercise}

\begin{exercise}[4]
Investigate profile-likelihood confidence intervals for nuisance-parameter
models.
\end{exercise}

\begin{exercise}[4]
Study BCa bootstrap confidence intervals.
\end{exercise}

\begin{exercise}[4]
Investigate bootstrap confidence intervals under dependence.
\end{exercise}

\begin{exercise}[4]
Study simultaneous confidence intervals using Bonferroni methods.
\end{exercise}

\begin{exercise}[4]
Investigate Scheff\'e and Tukey simultaneous inference procedures.
\end{exercise}

\begin{exercise}[4]
Study confidence ellipsoids for multivariate normal parameters.
\end{exercise}

\begin{exercise}[4]
Investigate robust sandwich standard errors and confidence intervals.
\end{exercise}

%%%%%%%%%%%%%%%%%%%%%%%%%%%%%%%%%%%%%%%%%%%%%%%%%%%%%%%%%%%%
\subsection{Conceptual Summary}
\label{subsec:confidence-interval-summary}
%%%%%%%%%%%%%%%%%%%%%%%%%%%%%%%%%%%%%%%%%%%%%%%%%%%%%%%%%%%%

A confidence interval is a random interval procedure designed to contain
an unknown parameter with a specified repeated-sampling probability.

Many intervals have the form

\[
\boxed{
\text{estimate}
\pm
\text{critical value}
\times
\text{standard error}.
}
\]

For a population mean with known \(\sigma\),

\[
\boxed{
\overline{X}
\pm
z^*
\frac{\sigma}{\sqrt{n}}.
}
\]

For a population mean with unknown \(\sigma\),

\[
\boxed{
\overline{X}
\pm
t^*_{n-1}
\frac{S}{\sqrt{n}}.
}
\]

For a large-sample population proportion,

\[
\boxed{
\widehat{p}
\pm
z^*
\sqrt{
\frac{
\widehat{p}(1-\widehat{p})
}{n}
}.
}
\]

For two independent means, Welch's interval uses

\[
\boxed{
(\overline{X}_1-\overline{X}_2)
\pm
t^*
\sqrt{
\frac{
S_1^2
}{
n_1
}
+
\frac{
S_2^2
}{
n_2
}
}.
}
\]

For paired data,

\[
\boxed{
\overline{D}
\pm
t^*
\frac{
S_D
}{
\sqrt{n}
}.
}
\]

For two independent proportions,

\[
\boxed{
(\widehat{p}_1-\widehat{p}_2)
\pm
z^*
\sqrt{
\frac{
\widehat{p}_1(1-\widehat{p}_1)
}{
n_1
}
+
\frac{
\widehat{p}_2(1-\widehat{p}_2)
}{
n_2
}
}.
}
\]

For a normal variance,

\[
\boxed{
\frac{
(n-1)S^2
}{
\sigma^2
}
\sim
\chi^2_{n-1}.
}
\]

Confidence intervals therefore combine

\[
\boxed{
\text{point estimation}
+
\text{sampling variability}
+
\text{probability}
}
\]

to quantify statistical uncertainty.

%%%%%%%%%%%%%%%%%%%%%%%%%%%%%%%%%%%%%%%%%%%%%%%%%%%%%%%%%%%%
\subsection{Section Connections}
\label{subsec:confidence-interval-connections}
%%%%%%%%%%%%%%%%%%%%%%%%%%%%%%%%%%%%%%%%%%%%%%%%%%%%%%%%%%%%

\chapterconnection{
Confidence Intervals combine the estimators of
Section~\ref{sec:statistical-estimation} with the sampling distributions
of Section~\ref{sec:sampling-distributions}. Exact normal, \(t\),
chi-square, and \(F\) pivots produce classical finite-sample intervals,
while asymptotic normality and bootstrap methods provide approximate
intervals for more general estimators. Confidence intervals communicate
both effect magnitude and statistical precision. The next section,
Hypothesis Testing, uses the same sampling distributions to evaluate
formal statistical claims. In many standard models, hypothesis tests and
confidence intervals are dual procedures: parameter values rejected by a
two-sided level-\(\alpha\) test are precisely those excluded from the
corresponding \(100(1-\alpha)\%\) confidence set.
}

%%%%%%%%%%%%%%%%%%%%%%%%%%%%%%%%%%%%%%%%%%%%%%%%%%%%%%%%%%%%
% SECTION SOURCE TRACKING
%%%%%%%%%%%%%%%%%%%%%%%%%%%%%%%%%%%%%%%%%%%%%%%%%%%%%%%%%%%%

%------------------------------------------------------------
% 8.5 Confidence Intervals
%------------------------------------------------------------
%
% Source basis:
%   - Standard mathematical statistics.
%   - Standard interval-estimation theory.
%   - Standard asymptotic and bootstrap inference.
%
% Used for:
%   - confidence-interval definitions;
%   - repeated-sampling interpretation;
%   - confidence coefficients;
%   - critical values;
%   - margin of error;
%   - known-variance mean intervals;
%   - one-sample t intervals;
%   - sample-size planning;
%   - proportion intervals;
%   - Wald interval limitations;
%   - Wilson score intervals;
%   - exact binomial intervals;
%   - two-sample mean intervals;
%   - Welch intervals;
%   - pooled intervals;
%   - paired intervals;
%   - difference-of-proportion intervals;
%   - variance intervals;
%   - standard-deviation intervals;
%   - ratio-of-variance intervals;
%   - one-sided confidence bounds;
%   - pivotal quantities;
%   - Wald intervals;
%   - likelihood-ratio intervals;
%   - score intervals;
%   - transformation-based intervals;
%   - delta-method intervals;
%   - bootstrap normal, percentile, and basic intervals;
%   - BCa preview;
%   - dependence-aware intervals;
%   - complex-survey intervals;
%   - confidence width and precision;
%   - prediction and tolerance interval distinctions;
%   - confidence-test duality;
%   - simultaneous intervals;
%   - confidence regions;
%   - applications and exercises.
%
% Notes:
%   - Hypothesis Testing is developed next in Section 8.6.
%   - Regression later introduces coefficient and prediction intervals.
%   - ANOVA and experimental design introduce simultaneous inference.
%   - Computational Statistics develops bootstrap methods more fully.
%
%%%%%%%%%%%%%%%%%%%%%%%%%%%%%%%%%%%%%%%%%%%%%%%%%%%%%%%%%%%%